\documentclass[conference]{IEEEtran}
\IEEEoverridecommandlockouts
\usepackage{cite}
\usepackage{amsmath,amssymb,amsfonts}
\usepackage{algorithmic}
\usepackage{graphicx}
\usepackage{textcomp}
\usepackage{xcolor}
\def\BibTeX{{\rm B\kern-.05em{\sc i\kern-.025em b}\kern-.08em
    T\kern-.1667em\lower.7ex\hbox{E}\kern-.125emX}}

\begin{document}

\title{PhishScope: An Explainable Hybrid Intelligence Framework for Static and Dynamic Phishing Website Analysis}

\author{
\IEEEauthorblockN{1\textsuperscript{st} Utkarsh Khandu Karale}
\IEEEauthorblockA{\textit{M.Tech. Computer Engineering} \\
\textit{Ashokrao Mane Group of Institutions}\\
Kolhapur, India \\
karaleutkarsh3696@gmail.com}
\and
\IEEEauthorblockN{2\textsuperscript{nd} Prof. Vikas Patil}
\IEEEauthorblockA{\textit{Department of Computer Engineering} \\
\textit{Ashokrao Mane Group of Institutions}\\
Kolhapur, India \\
guide.email@example.com}
}

\maketitle

\begin{abstract}
Phishing attacks have evolved into highly sophisticated threats utilizing dynamic web content and multi-stage redirections. Traditional detection systems relying solely on static features are increasingly inadequate and often operate as uninterpretable ``black boxes.'' In this paper, we propose \textit{PhishScope}, an explainable hybrid intelligence framework that integrates multi-source static analysis (URL, HTML, Domain, SSL) with isolated browser-based dynamic behavioral execution. We employ a hybrid machine learning classification approach utilizing Random Forest and XGBoost. To establish trust and transparency, we integrate Explainable AI (XAI) techniques, specifically SHAP and LIME, to generate evidence-based risk assessments. Experimental results demonstrate that PhishScope achieves high detection accuracy while providing human-interpretable explanations for its predictions.
\end{abstract}

\begin{IEEEkeywords}
Phishing Detection, Explainable AI, SHAP, LIME, Hybrid Machine Learning, Cybersecurity
\end{IEEEkeywords}

\section{Introduction}
% Introduce the problem of phishing, the limitations of current static/black-box models, and outline the contributions of PhishScope.
Phishing remains one of the most prominent cybersecurity threats...

\section{Related Work}
% Summarize the 12 papers from the synopsis literature review here.
Recent advancements in phishing detection have heavily utilized Machine Learning...

\section{Proposed Methodology}
% Describe Data Collection, Static Extractor, Dynamic Sandbox, and the ML Pipeline.
\subsection{Multi-Source Static Feature Extraction}
We extract features across four dimensions: URL, HTML, Domain, and SSL...

\subsection{Dynamic Behavioral Sandbox}
To capture runtime execution patterns...

\subsection{Hybrid ML Classification}
We propose an ensemble of Random Forest and XGBoost...

\section{Explainable AI (XAI) Integration}
% Explain how SHAP and LIME are used for Risk Scoring.
To demystify the ``black-box'' nature of our classifier...

\section{Experimental Setup and Results}
% Detail datasets, metrics, and display results/charts.
\subsection{Dataset}
\subsection{Performance Metrics}
\subsection{Results Analysis}

\section{Conclusion and Future Scope}
In this paper, we presented PhishScope...

\bibliographystyle{IEEEtran}
\bibliography{references}

\end{document}
